\documentclass[aspectratio=169,10pt,t]{beamer}
\usepackage{ac1}
\aulainfo{23}{Arquivos}{Leitura e escrita de arquivos}

\begin{document}
\capa

\begin{frame}{Roteiro}
  \begin{itemize}
    \item Por que arquivos? Memória volátil \emph{versus} armazenamento permanente
    \item Abrindo, lendo e fechando (Cap.~6.4)
    \item Escrevendo em arquivos (6.5)
    \item O comando \texttt{with} e tratamento de erros
    \item Arquivos CSV: dados em tabela
  \end{itemize}
\end{frame}

\begin{frame}{Programas e dados persistentes}
  \begin{columns}[T]
    \column{0.55\textwidth}
    \begin{itemize}
      \item As variáveis vivem na \textbf{RAM}: desaparecem quando o programa termina
      \item Para guardar resultados, ler medições ou trocar dados com outros programas usamos \textbf{arquivos}
      \item Em Python, um arquivo de texto é visto como uma \emph{sequência de linhas} (strings)
    \end{itemize}
    \column{0.42\textwidth}
    \begin{exampleblock}{Exemplos}
      \begin{itemize}\small
        \item resultados de um ensaio
        \item log de um sensor
        \item lista de alunos
        \item planilhas exportadas em CSV
      \end{itemize}
    \end{exampleblock}
  \end{columns}
\end{frame}

\begin{frame}[fragile]{Abrindo e lendo: \texttt{open}}
  Arquivo \texttt{nomes.txt}:
\begin{lstlisting}[style=saida]
João Araujo Ribeiro
Grace Hopper
Ada Lovelace
\end{lstlisting}
\begin{lstlisting}
arquivo = open('nomes.txt', 'r')     # 'r' = read (leitura)
linhas = arquivo.readlines()         # lista de strings
arquivo.close()                      # libera o arquivo
for linha in linhas:
    print(linha)
\end{lstlisting}
  \begin{alertblock}{Por que as linhas saem com uma linha em branco a mais?}
    Cada linha lida termina com \texttt{\textbackslash n}, e \texttt{print} acrescenta outro. Resolva com \texttt{rstrip()}.
  \end{alertblock}
  \fonte{code/Cap6/prog.6.12\_abrearquivo.py}
\end{frame}

\begin{frame}[fragile]{Formas de ler}
\begin{lstlisting}
# 1) linha por linha (não carrega o arquivo todo)
arquivo = open('nomes.txt', 'r')
for linha in arquivo:
    print(linha.rstrip())
arquivo.close()

# 2) tudo como uma lista, já sem o \n
arquivo = open('nomes.txt', 'r')
linhas = [linha.rstrip() for linha in arquivo.readlines()]
arquivo.close()

# 3) tudo em uma única string
texto = open('nomes.txt').read()
linhas = texto.split('\n')
\end{lstlisting}
  \fonte{code/Cap6/prog.6.13, prog.6.14, prog.6.15 (abre arquivo)}
\end{frame}

\begin{frame}[fragile]{Escrevendo: modos \texttt{'w'} e \texttt{'a'}}
\begin{lstlisting}[basicstyle=\ttfamily\scriptsize]
import math

arquivo = open('tabelaSenoCosseno.txt', 'w')    # 'w' = write
arquivo.write(f"{'Graus':>9}{'Radianos':>9}{'Seno':>9}{'Cosseno':>9}\n")
for graus in range(0, 361, 30):
    rad = math.radians(graus)
    arquivo.write(f"{graus:>9.2f}{rad:>9.2f}"
                  f"{math.sin(rad):>9.2f}{math.cos(rad):>9.2f}\n")
arquivo.close()
\end{lstlisting}
  \begin{center}\small
  \begin{tabular}{cl}
    \toprule
    Modo & Efeito\\
    \midrule
    \texttt{'r'} & leitura (padrão); erro se o arquivo não existir\\
    \texttt{'w'} & escrita; \alert{apaga} o conteúdo anterior!\\
    \texttt{'a'} & escrita no final (\emph{append}), mantendo o conteúdo\\
    \bottomrule
  \end{tabular}
  \end{center}
  \fonte{code/Cap6/prog.6.17\_tabela\_senos\_arquivo\_f.py}
\end{frame}

\begin{frame}[fragile]{\texttt{write} não acrescenta \texttt{\textbackslash n}}
\begin{lstlisting}
arq = open('saida.txt', 'w')
arq.write('linha 1')
arq.write('linha 2')        # saiu "linha 1linha 2" !
arq.write('linha 3\n')      # \n explícito
print('linha 4', file=arq)  # print grava com \n automático
arq.close()
\end{lstlisting}
  \begin{itemize}
    \item \texttt{write} só aceita \alert{strings}: converta números com \texttt{str()} ou f-string
    \item Esquecer de fechar (\texttt{close}) pode deixar dados sem gravar
  \end{itemize}
\end{frame}

\begin{frame}[fragile]{O comando \texttt{with}}
\begin{lstlisting}
with open('nomes.txt', 'r') as arquivo:
    for linha in arquivo:
        print(linha.rstrip())
# aqui o arquivo já foi fechado automaticamente

with open('copia.txt', 'w') as saida:
    with open('nomes.txt', 'r') as entrada:
        for linha in entrada:
            saida.write(linha)
\end{lstlisting}
  \begin{itemize}
    \item \texttt{with} fecha o arquivo \alert{sempre}, mesmo que ocorra um erro: é a forma recomendada
    \item O segundo exemplo é o Exercício 6.14 (cópia de arquivo)
  \end{itemize}
\end{frame}

\begin{frame}[fragile]{Arquivos que não existem}
\begin{lstlisting}
nome = input('Nome do arquivo: ')
try:
    with open(nome, 'r') as f:
        conteudo = f.read()
except FileNotFoundError:
    print('Arquivo não encontrado:', nome)
except PermissionError:
    print('Sem permissão para ler', nome)
else:
    print(len(conteudo), 'caracteres lidos')
\end{lstlisting}
  \begin{itemize}
    \item Combina o que vimos na Aula 13 (\texttt{try}/\texttt{except}) com arquivos
    \item O caminho é relativo à pasta onde o programa \emph{executa} (não à do arquivo \texttt{.py})
  \end{itemize}
\end{frame}

\begin{frame}[fragile]{Exemplo: processando um arquivo de texto}
  Converter o conteúdo de um arquivo para maiúsculas e gravar o resultado (Exercício 6.10):
\begin{lstlisting}
with open('nomes.txt', 'r') as entrada:
    texto = entrada.read()
with open('nomes_maiusculas.txt', 'w') as saida:
    saida.write(texto.upper())
\end{lstlisting}
  \textbf{Variações (Exercícios 6.11 e 6.12):}
  \begin{itemize}
    \item \texttt{texto.lower()} para minúsculas
    \item Primeira letra de cada palavra maiúscula: \texttt{texto.title()} ou processar \texttt{split()} palavra a palavra
  \end{itemize}
\end{frame}

\begin{frame}[fragile]{Dados em tabela: arquivos CSV}
  CSV = \emph{Comma-Separated Values}: uma linha por registro, campos separados por vírgula.
\begin{lstlisting}[style=saida]
id_peca,maquina,operador,turno,tempo_ciclo_seg,status
1000,M1,Ana,Noite,31.6,Aprovado
1001,M1,Elis,Tarde,34.5,Aprovado
1002,M2,Bruno,Noite,24.9,Descarte
\end{lstlisting}
\begin{lstlisting}
import csv
with open('dadosproducao.csv', newline='', encoding='utf-8') as f:
    leitor = csv.DictReader(f)          # cada linha vira um dicionário
    registros = list(leitor)
print(registros[0]['maquina'], float(registros[0]['tempo_ciclo_seg']))
\end{lstlisting}
  \small Dica: \texttt{DictReader} + o que vimos de dicionários (Aula 20) resolve muitos problemas. O módulo \texttt{csv} cuida de aspas e vírgulas dentro dos campos.\\
  \fonte{code/Projeto1/ler\_producao\_dict.py}
\end{frame}

\begin{frame}[fragile]{Gravando CSV}
\begin{lstlisting}
import csv
ensaios = [(5.0, 0.25), (10.0, 0.52), (15.0, 0.76)]
with open('ensaios.csv', 'w', newline='', encoding='utf-8') as f:
    escritor = csv.writer(f)
    escritor.writerow(['tensao_V', 'corrente_A'])
    escritor.writerows(ensaios)
\end{lstlisting}
  \begin{itemize}
    \item Todo programa de planilha (Excel, LibreOffice) abre CSV
    \item Na parte final do curso, o \textbf{Pandas} lê CSV em uma única linha: \texttt{pd.read\_csv('ensaios.csv')}
  \end{itemize}
\end{frame}

\begin{frame}{Atividades}
  \begin{enumerate}
    \item Escreva um programa que leia \texttt{nomes.txt} e grave em outro arquivo os nomes em ordem alfabética, um por linha.
    \item Leia um arquivo de notas (uma por linha) e escreva a média, a maior e a menor em um arquivo \texttt{resumo.txt}.
    \item Grave em \texttt{tabela.txt} a tabela Celsius \(\to\) Fahrenheit de 0 a 100 de 10 em 10 (Exercício 6.13).
    \item Leia \texttt{dadosproducao.csv} e conte quantas peças cada máquina produziu (dicionário de contagens, Aula 20).
  \end{enumerate}
\end{frame}

\exercicios{6}{%
  \begin{itemize}\small
    \item Leitura: \textbf{6.10, 6.11, 6.12} (conversão de maiúsculas/minúsculas)
    \item Escrita: \textbf{6.13} (tabela em arquivo), \textbf{6.14} (cópia de arquivo), \textbf{6.15} (pontos de \(\pi\) em arquivo)
  \end{itemize}}

\resumo{
  \item \texttt{open(nome, modo)}: \texttt{'r'}, \texttt{'w'} (apaga!) e \texttt{'a'}
  \item Arquivo de texto é percorrido linha a linha; \texttt{rstrip()} tira o \texttt{\textbackslash n}
  \item \texttt{with} fecha o arquivo automaticamente; trate \texttt{FileNotFoundError}
  \item CSV: use o módulo \texttt{csv} (ou Pandas, nas aulas finais)
  \item Próxima aula: revisão, estilo e boas práticas
}

\end{document}
